% GENERATED FILE. Edit canonical lessons, then run npm run generate:books.
% canonical-source-hash: fnv1a64:b97298a0892750ad
% canonical-lessons: GE-C11-brot, GE-C11-brot-bread, GE-C11-wasser, GE-C11-wasser-water, GE-C11-wein, GE-C11-wein-vinum, GE-C11-wasser-wein

\chapter{Bread, Water, Wine}
\label{ch:food}
\hlchaptermodality{17}

\begin{chapteropening}
\textbf{By the end of this chapter:} \emph{I can name and request bread, water, and wine in German.}

Three things you could put on one table and two histories between them. Brot is Germanic braudam, English bread, where French took Latin panis and German did not. Wasser is Germanic watar with German's own later shift turning the t into ss, the same shift that put the tt into Mutter. Wein is Latin vinum, borrowed whole with the grapevine the legions carried north -- and you can tell it is borrowed without being told, because the inherited pair have drifted from their English twins and the loan has not.
\end{chapteropening}

\section[das Brot]{das Brot — \textquotedblleft{}bread\textquotedblright{}}

\label{lesson:GE-C11-brot}



\begin{quote}
\emph{Der Vater}. \emph{Die Mutter}. You were given a third article in the first chapter and have had no noun for it. Here it is.
\end{quote}

\subsection*{You'll want to know: das Brot}
\begin{quote}
\textbf{das Brot} — \textquotedblleft{}bread.\textquotedblright{}
\end{quote}

A \emph{\textbf{das}}-word. That completes the set on real nouns: \emph{der Vater}, \emph{die Mutter}, \emph{das Brot} — one of each, and now you have met all three doing their job rather than sitting in a list.

\begin{sounds}
\begin{itemize}
\raggedright
  \item \texttt{long-o} — a long \emph{o}, the vowel of English \emph{boat}, held steady.
  \item \texttt{r-uvular} — the \emph{r} is made far back in the throat.
  \item \texttt{d-t-final} — and the final \emph{t} is crisp.
\end{itemize}

\emph{broht}. One syllable, and capitalised on the page like every German noun, whether it names a person or a loaf.
\end{sounds}

\subsection*{Your turn}
Say these aloud:

\begin{itemize}
\raggedright
  \item \textquotedblleft{}das Brot\textquotedblright{} — broht
  \item one noun of each gender — \textquotedblleft{}der Vater, die Mutter, das Brot\textquotedblright{}
  \item the three articles alone — \textquotedblleft{}der, die, das\textquotedblright{}
  \item why it has a capital — \textquotedblleft{}every German noun does\textquotedblright{}
\end{itemize}

\emph{Twice through:} \textquotedblleft{}das Brot.\textquotedblright{}

\begin{tcolorbox}[breakable,colback=teal!4,colframe=teal!35!black,title={Before you move on}]
Say \textquotedblleft{}bread.\textquotedblright{} (\textbf{Das Brot}.) Which article? (\emph{\textbf{Das}} — the neuter.) Say all three with a noun each. (\textbf{Der Vater, die Mutter, das Brot}.) Why the capital \emph{B}? (\textbf{German capitalises every noun}.) Next: where the word comes from, and where it does not.
\end{tcolorbox}
\section[Brot und bread]{Brot and bread}

\label{lesson:GE-C11-brot-bread}



\begin{quote}
A word can be interesting for what happened to it. This one is interesting for what did not.
\end{quote}

\begin{cousinweb}
\emph{Brot} comes from Germanic \emph{*braudam}, and English \textbf{bread} comes from the same word. No Latin anywhere in it.

Now the other side of the line. Latin \emph{pānis} went into every Romance language and stayed there:

\noindent
\begin{tabularx}{\linewidth}{@{}>{\raggedright\arraybackslash}X>{\raggedright\arraybackslash}X>{\raggedright\arraybackslash}X@{}}
\toprule
\textbf{language} & \textbf{the word} & \textbf{from} \\
\midrule
German & \emph{Brot} & Germanic \emph{*braudam} \\
French & \emph{pain} & Latin \emph{pānis} \\
Italian & \emph{pane} & Latin \emph{pānis} \\
\bottomrule
\end{tabularx}
\end{cousinweb}

\begin{culture}
Rome sold German a calendar and a clock. It did not sell German bread.

The pattern is the one you have been given four times now and it holds again here: a thing people handle every single day already has a word, and an incoming language does not get to replace it. Months were an administrative object and arrived with their labels. A loaf on a table was not waiting for a name.

That is why \emph{Brot} sits with \emph{Vater} and \emph{Mutter} rather than with \emph{Januar}.
\end{culture}

\subsection*{Your turn}
Say these aloud:

\begin{itemize}
\raggedright
  \item the Germanic pair — \textquotedblleft{}Brot, bread\textquotedblright{}
  \item the Romance word it is not — \textquotedblleft{}pain\textquotedblright{}
  \item which kind of word this is — \textquotedblleft{}inherited, not bought\textquotedblright{}
\end{itemize}

\emph{Twice through:} \textquotedblleft{}Brot, bread.\textquotedblright{}

\begin{tcolorbox}[breakable,colback=teal!4,colframe=teal!35!black,title={Before you move on}]
What is the English twin of \emph{Brot}? (\emph{\textbf{Bread}}.) What does French say? (\emph{\textbf{Pain}}, \textbf{from Latin} \emph{pānis}.) Did German borrow that? (\textbf{No}.) Why not? (\textbf{A daily thing already had a word}.) Next: the second of the three.
\end{tcolorbox}
\section[das Wasser]{das Wasser — \textquotedblleft{}water\textquotedblright{}}

\label{lesson:GE-C11-wasser}



\begin{quote}
Bread, and then something to drink beside it. Its first letter you have met twice already.
\end{quote}

\subsection*{You'll want to know: das Wasser}
\begin{quote}
\textbf{das Wasser} — \textquotedblleft{}water.\textquotedblright{}
\end{quote}

Another \emph{\textbf{das}}-word, so it takes the same article as \emph{Brot}. Two neuters in a row is a useful accident: it lets you say \emph{das Brot und das Wasser} without having to decide anything.

\begin{sounds}
\begin{itemize}
\raggedright
  \item \texttt{w-as-v} — German \emph{w} is an English \textbf{v}. \emph{Wasser} opens \emph{VAH-}.
  \item \texttt{vowel-a-german} — a clean open \emph{a}.
  \item \texttt{r-final} — and the soft closing \emph{r}.
\end{itemize}

\emph{VAH-ser}. You have had that \emph{w} on \emph{Winter}, and inside the \emph{shv-} of \emph{Schwester}; this is the third time and the first where it sits on its own at the front of a word.
\end{sounds}

\subsection*{Your turn}
Say these aloud:

\begin{itemize}
\raggedright
  \item \textquotedblleft{}das Wasser\textquotedblright{} — VAH-ser
  \item the letter and the sound — \textquotedblleft{}W on the page, v in the mouth\textquotedblright{}
  \item the two together — \textquotedblleft{}das Brot und das Wasser\textquotedblright{}
\end{itemize}

\emph{Twice through:} \textquotedblleft{}das Wasser.\textquotedblright{}

\begin{tcolorbox}[breakable,colback=teal!4,colframe=teal!35!black,title={Before you move on}]
Say \textquotedblleft{}water.\textquotedblright{} (\textbf{Das Wasser}.) Which article? (\emph{\textbf{Das}}, like \emph{Brot}.) How is the \emph{W} said? (\textbf{As a} \emph{v}.) Where did you meet that before? (\textbf{On} \emph{Winter}, \textbf{and inside} \emph{Schwester}.) Next: its English twin.
\end{tcolorbox}
\section[Wasser und water]{Wasser and water}

\label{lesson:GE-C11-wasser-water}



\begin{quote}
Say \emph{Wasser} and then \textbf{water}. The difference is one consonant, and you have met it before.
\end{quote}

\begin{cousinweb}
Both come from Germanic \emph{*watar}. English kept the \textbf{t}. German turned it into \textbf{ss}.

That is German's own later shift — the second of the two you met on \emph{Mutter}, where English stopped at \emph{mother} and German kept going. It ran through the vocabulary the same way, so once you have it you can predict it:

\noindent
\begin{tabularx}{\linewidth}{@{}>{\raggedright\arraybackslash}X>{\raggedright\arraybackslash}X@{}}
\toprule
\textbf{English} & \textbf{German} \\
\midrule
wa\textbf{t}er & Wa\textbf{ss}er \\
mo\textbf{th}er & Mu\textbf{tt}er \\
\bottomrule
\end{tabularx}
\end{cousinweb}

\begin{culture}
Two chapters, four inherited words, and the same two shifts explaining every difference between the German and the English.

That is the payoff of a sound law: you stop memorising pairs. \emph{Wasser} is not a word to be learned beside \emph{water} — it \textbf{is} \emph{water}, with one predictable change applied, and the change is one you can now run yourself.

And like \emph{Brot}, nothing here was bought. Water is as everyday as a word gets.
\end{culture}

\subsection*{Your turn}
Say these aloud:

\begin{itemize}
\raggedright
  \item the pair — \textquotedblleft{}Wasser, water\textquotedblright{}
  \item what changed — \textquotedblleft{}t became ss\textquotedblright{}
  \item the other pair with the same shift — \textquotedblleft{}Mutter, mother\textquotedblright{}
\end{itemize}

\emph{Twice through:} \textquotedblleft{}Wasser, water.\textquotedblright{}

\begin{tcolorbox}[breakable,colback=teal!4,colframe=teal!35!black,title={Before you move on}]
What is the English twin? (\emph{\textbf{Water}}.) What did German do to the \emph{t}? (\textbf{Made it} \emph{ss}.) Which other word did the same shift? (\emph{\textbf{Mutter}}.) Inherited or borrowed? (\textbf{Inherited}.) Next: the one on this table that was borrowed.
\end{tcolorbox}
\section[der Wein]{der Wein — \textquotedblleft{}wine\textquotedblright{}}

\label{lesson:GE-C11-wein}



\begin{quote}
Two \emph{das}-words so far. The third drink takes a different article, and that is the smaller of the two things different about it.
\end{quote}

\subsection*{You'll want to know: der Wein}
\begin{quote}
\textbf{der Wein} — \textquotedblleft{}wine.\textquotedblright{}
\end{quote}

A \emph{\textbf{der}}-word, where \emph{Brot} and \emph{Wasser} are both \emph{das}. So the three now run one of each: \emph{das Brot}, \emph{das Wasser}, \emph{der Wein}.

\begin{sounds}
\begin{itemize}
\raggedright
  \item \texttt{w-as-v} — the \emph{W} is a \textbf{v} again, as in \emph{Wasser}.
  \item \texttt{diphthong-ei} — and \emph{ei} is the vowel of English \textbf{eye}, which you were given on \emph{eins}.
\end{itemize}

Put the two together and \emph{Wein} comes out sounding like English \textbf{vine}. Hold on to that, because the next lesson is about why.
\end{sounds}

\subsection*{Your turn}
Say these aloud:

\begin{itemize}
\raggedright
  \item \textquotedblleft{}der Wein\textquotedblright{} — like English \textquotedblleft{}vine\textquotedblright{}
  \item the two vowels you know — \textquotedblleft{}ei is eye, as in eins\textquotedblright{}
  \item all three with their articles — \textquotedblleft{}das Brot, das Wasser, der Wein\textquotedblright{}
\end{itemize}

\emph{Twice through:} \textquotedblleft{}der Wein.\textquotedblright{}

\begin{tcolorbox}[breakable,colback=teal!4,colframe=teal!35!black,title={Before you move on}]
Say \textquotedblleft{}wine.\textquotedblright{} (\textbf{Der Wein}.) Which article? (\emph{\textbf{Der}} — not \emph{das}, like the other two.) What does \emph{ei} sound like? (\textbf{English} \emph{eye}.) So what English word does it sound like? (\emph{\textbf{Vine}}.) Next: why it does.
\end{tcolorbox}
\section[Wein und vinum]{Wein and vinum}

\label{lesson:GE-C11-wein-vinum}



\begin{quote}
\emph{Wasser} and \emph{water} differ by a consonant. \emph{Wein} and \textbf{wine} differ by nothing you can hear. That difference is the lesson.
\end{quote}

\begin{cousinweb}
\emph{Wein} is \textbf{not} an inherited Germanic word. It was \textbf{borrowed from Latin \emph{vīnum}} in Roman times, when the legions brought the grapevine north of the Alps.

Here is how you can tell without being told. Inherited words drift apart, because a sound law runs through them:

\noindent
\begin{tabularx}{\linewidth}{@{}>{\raggedright\arraybackslash}X>{\raggedright\arraybackslash}X@{}}
\toprule
\textbf{pair} & \textbf{how far apart} \\
\midrule
\emph{Wasser} / water & a whole consonant — \emph{ss} against \emph{t} \\
\emph{Mutter} / mother & two shifts of distance \\
\emph{Wein} / wine / \emph{vīnum} & \textbf{no distance at all} \\
\bottomrule
\end{tabularx}

Three words that match this closely are not three cousins who happen to look alike. They are \textbf{one word, borrowed}, arriving too late for the law to touch it.
\end{cousinweb}

\begin{culture}
The vine did not grow north of the Alps until Rome carried it there. A people who have never had a thing have no word for it, and when the thing arrives the word arrives with it — which is the same reason \emph{Uhr} is Latin and the twelve months are Latin.

So one table holds both halves of the rule. \emph{Brot} and \emph{Wasser} were always here and are German's own. \emph{Wein} came on a cart, and its name was in the same load.
\end{culture}

\subsection*{Your turn}
Say these aloud:

\begin{itemize}
\raggedright
  \item the three that match — \textquotedblleft{}Wein, wine, vinum\textquotedblright{}
  \item what that closeness means — \textquotedblleft{}one loan, not three cousins\textquotedblright{}
  \item what came with the word — \textquotedblleft{}the vine itself\textquotedblright{}
  \item the two that were always here — \textquotedblleft{}das Brot, das Wasser\textquotedblright{}
\end{itemize}

\emph{Twice through:} \textquotedblleft{}Wein, wine, vinum — one loan.\textquotedblright{}

\begin{tcolorbox}[breakable,colback=teal!4,colframe=teal!35!black,title={Before you move on}]
Is \emph{Wein} inherited? (\textbf{No} — \textbf{borrowed from Latin} \emph{vīnum}.) How can you tell? (\textbf{It is too close} — a law would have moved it.) What arrived with the word? (\textbf{The grapevine}.) And \emph{Wasser}? (\textbf{Always here}.)
\end{tcolorbox}
\section[Practice]{Practice — one table, two histories}

\label{lesson:GE-C11-wasser-wein}



\begin{quote}
No new words. Three things you could put on a table, and one of them is a different kind of word from the other two.
\end{quote}

\begin{grammarlens}[title={Grammar Lens: three nouns, three articles}]
\noindent
\begin{tabularx}{\linewidth}{@{}>{\raggedright\arraybackslash}X>{\raggedright\arraybackslash}X@{}}
\toprule
\textbf{German} & \textbf{English} \\
\midrule
\emph{das Brot} & the bread \\
\emph{das Wasser} & the water \\
\emph{der Wein} & the wine \\
\bottomrule
\end{tabularx}

Two neuters and a masculine, every one of them capitalised. There is no rule behind which is which — the article comes with the noun and is learned with it, which is why every word in this book arrives wearing one.
\end{grammarlens}

\begin{culture}
Now the histories, which do have a rule behind them:

\begin{itemize}
\raggedright
  \item \emph{\textbf{Brot}} — Germanic \emph{*braudam}, English \textbf{bread}. French took Latin \emph{pānis} and German did not.
  \item \emph{\textbf{Wasser}} — Germanic \emph{*watar}, English \textbf{water}, with German's own later shift turning the \emph{t} into \emph{ss}.
  \item \emph{\textbf{Wein}} — Latin \emph{vīnum}, borrowed whole, matching English \textbf{wine} and the Latin too closely to be anything else.
\end{itemize}

Two were always here; one came north on a cart with the plant it names. You can read that off the words themselves without knowing any history: the inherited pair have \textbf{drifted} from their English twins, and the borrowed one has not.
\end{culture}

\subsection*{Your turn}
Say these aloud:

\begin{itemize}
\raggedright
  \item all three with their articles — \textquotedblleft{}das Brot, das Wasser, der Wein\textquotedblright{}
  \item the two with a v sound — \textquotedblleft{}Wasser, Wein\textquotedblright{}
  \item which two are German's own — \textquotedblleft{}Brot, Wasser\textquotedblright{}
  \item which one Rome supplied — \textquotedblleft{}Wein\textquotedblright{}
  \item how you can tell — \textquotedblleft{}it did not drift\textquotedblright{}
\end{itemize}

\emph{Twice through:} \textquotedblleft{}das Brot, das Wasser, der Wein.\textquotedblright{}

\begin{tcolorbox}[breakable,colback=teal!4,colframe=teal!35!black,title={Before you move on}]
Say the three with their articles. (\textbf{Das Brot, das Wasser, der Wein}.) Which is masculine? (\emph{\textbf{Der Wein}}.) Which two did German inherit? (\emph{\textbf{Brot}} \textbf{and} \emph{\textbf{Wasser}}.) Which came from Latin? (\emph{\textbf{Wein}}, \textbf{from} \emph{vīnum}.) And what arrived with it? (\textbf{The vine}.)
\end{tcolorbox}
